\section*{ВВЕДЕНИЕ}
\addcontentsline{toc}{section}{ВВЕДЕНИЕ}

Веб-платформы для поиска работы стали ключевым инструментом на современном рынке труда, обеспечивая удобное взаимодействие между соискателями и работодателями. Они автоматизируют процессы подбора персонала, упрощая создание резюме, публикацию вакансий, поиск по критериям и отправку откликов. Благодаря развитию информационных технологий такие платформы предлагают интуитивный интерфейс, высокую производительность и доступность для миллионов пользователей. Они заменили традиционные способы поиска работы, такие как объявления в газетах, и стали стандартом для эффективного найма.

Разработка веб-платформы для поиска работы предполагает создание системы, которая объединяет функциональность для регистрации пользователей, управления резюме и вакансиями, поиска и обработки откликов. Такие платформы используют современные веб-технологии, базы данных и системы управления для обеспечения надежности, безопасности и удобства.

Веб-платформы помогают соискателям находить подходящие вакансии, а работодателям — подбирать кандидатов, соответствующих их требованиям. Они обеспечивают быстрый доступ к актуальной информации, упрощают коммуникацию и повышают эффективность рынка труда.

\emph{Цель настоящей работы} – разработка веб-платформы для поиска работы, которая облегчает взаимодействие соискателей и работодателей, упрощает подбор персонала и предоставляет удобный доступ к информации о вакансиях и резюме. Для достижения цели необходимо решить \emph{следующие задачи:}
\begin{itemize}
\item провести анализ предметной области веб-платформ для поиска работы;
\item разработать концептуальную модель платформы;
\item спроектировать платформу с учетом функциональных и нефункциональных требований;
\item реализовать платформу с использованием современных веб-технологий.
\end{itemize}

\emph{Структура и объем работы.} Отчет состоит из введения, 4 разделов основной части, заключения, списка использованных источников, 2 приложений. Текст выпускной квалификационной работы равен \formbytotal{lastpage}{страниц}{е}{ам}{ам}.

\emph{Во введении} сформулирована цель работы, определены задачи разработки, описана структура работы и краткое содержание разделов.

\emph{В первом разделе} анализируется предметная область: принципы работы веб-платформ для поиска работы, их функциональность, преимущества и недостатки для пользователей.

\emph{Во втором разделе} представлено техническое задание, включая требования к функциональности платформы, пользовательскому интерфейсу и оформлению документации.

\emph{В третьем разделе} описаны проектные решения: архитектура платформы, структура базы данных и макеты интерфейса.

\emph{В четвертом разделе} приведены реализованные модули, классы и методы, а также результаты тестирования платформы.

В заключении подводятся итоги разработки, оцениваются достигнутые результаты и предлагаются направления для дальнейшего развития.

В приложении А представлен графический материал.
В приложении Б представлены фрагменты исходного кода. 
